\begin{definition}[Minimal Iterator Relation]
\label{def:minimal-iterator-relation}
Let $(P,S,1)$ be a Peano system, and let $(W,c,g)$ be iterator data for it.
The \emph{minimal iterator relation} for $(W,c,g)$ is
\[
R^*
=
\bigcap\{R\subseteq P\times W : R \text{ is an iterator relation for }
(W,c,g)\}.
\]
\end{definition}
\end{definitionbox}

\begin{formalizationrecord}{Lean 4 Verification Record}
\FormalizationSource{Feferman, \emph{The Number Systems}, Section 3.4; Mendelson, \emph{Number Systems and the Foundations of Analysis}, Section 2.2}
\Formalizes{def:minimal-iterator-relation}
\textbf{Lean module:} \texttt{LRA.VolumeII.PeanoSystems.Recursion}.\par
\LeanDeclaration{minimal\_iterator\_relation}
\textbf{Status:} checked against \texttt{lra-lean}.\par

\begin{leancode}
def minimal_iterator_relation
    (ps : PeanoSystem)
    (data : IteratorData ps)
    (element : ps.carrier)
    (value : data.target) : Prop :=
  forall relation : ps.carrier -> data.target -> Prop,
    iterator_relation ps data relation ->
    relation element value
\end{leancode}
\end{formalizationrecord}

\begin{remark*}[Standard quantified statement]
\[
(n,w)\in R^*
\Longleftrightarrow
\forall R\subseteq P\times W\;
\bigl(R \text{ is an iterator relation for } (W,c,g)
\Longrightarrow (n,w)\in R\bigr).
\]
\end{remark*}

\begin{remark*}[Predicate reading]
Let $\mathcal C_{\mathrm{it}}$ be the closure condition
\[
\mathcal C_{\mathrm{it}}(R)
\Longleftrightarrow
\operatorname{IteratorRelation}(R,(P,S,1),W,c,g).
\]
Then the definition says
\[
\operatorname{MinimalRelation}(R^*,\mathcal C_{\mathrm{it}}).
\]
\end{remark*}

\begin{remark*}[Negated quantified statement]
\[
(n,w)\notin R^*
\Longleftrightarrow
\exists R\subseteq P\times W\;
\bigl(R \text{ is an iterator relation for } (W,c,g)\land (n,w)\notin R\bigr).
\]
\end{remark*}

\begin{remark*}[Failure modes]
\begin{description}
  \item[Exposition.]
  A pair fails to belong to the minimal iterator relation exactly when some
  iterator relation omits that pair. Such a pair is not forced by the base and
  successor-closure clauses.
  \item[Omitted pair.]
  \[
  \neg (n,w)\in R^*.
  \]
  \[
  (n,w)\notin R^*
  \Longleftrightarrow
  \exists R\subseteq P\times W\;
  \bigl(R \text{ is an iterator relation for } (W,c,g)\land (n,w)\notin R\bigr).
  \]
\end{description}
\end{remark*}

\begin{remark*}[Interpretation]
The minimal iterator relation consists exactly of the stage-value pairs forced
by the base pair and the successor closure rule. This is the $R^*$ construction
used in Feferman's proof of recursive definition, and in the related
set-theoretic presentations cited here
\citep{FefermanNumberSystems1964,hrbacek_jech1999,chi_recursion_wustl}.
\end{remark*}

\begin{dependencies}
\begin{itemize}
  \item \hyperref[def:iterator-relation]{Iterator Relation}
\end{itemize}
\end{dependencies}

